\label{sec:method}

\subsection{Model Selection}

We select Random Forest \cite{breiman2001} as the primary classifier on the following
grounds:

\begin{enumerate}
  \item \textbf{Heterogeneous feature handling}: the 24-dimensional feature vector combines
        integer counts, continuous statistics, and binary flags. Random Forest requires no
        feature normalisation and is invariant to monotonic feature transformations.

  \item \textbf{Class imbalance robustness}: the \texttt{class\_weight=balanced} setting
        assigns per-sample weights inversely proportional to class frequencies, effectively
        oversampling the malicious minority without synthetic data augmentation.

  \item \textbf{Interpretability}: feature importance scores are derived directly from
        mean decrease in Gini impurity across all trees, providing an inherently
        interpretable ranking aligned with the ablation study in Section~\ref{sec:exp}.

  \item \textbf{Low latency}: prediction requires $O(D \cdot T)$ operations for $D$
        feature dimensions and $T$ trees, yielding sub-millisecond per-session inference on
        a single CPU core.
\end{enumerate}

\subsection{Training Protocol}

Hyperparameters are selected via \texttt{RandomizedSearchCV} (20 random configurations,
evaluation metric: macro F1) over the following grids:

\begin{itemize}
  \item \texttt{n\_estimators}: $\{50, 100, 200\}$
  \item \texttt{max\_depth}: $\{5, 10, \text{None}\}$
  \item \texttt{min\_samples\_split}: $\{2, 5, 10\}$
  \item \texttt{max\_features}: $\{\text{sqrt}, \text{log2}\}$
\end{itemize}

Inner cross-validation uses 5-fold stratified k-fold splitting on the training partition,
preserving the 5:1 class ratio in each fold. The best configuration is refitted on the
combined train+validation set, and final metrics are reported on the held-out test
partition.

\subsection{Evaluation Metrics}

We report accuracy, precision, recall, macro F1, per-class F1, and ROC-AUC. For the
ablation study, the primary metric is macro F1, which weights benign and malicious classes
equally, penalising both false positives (operational disruption) and false negatives
(missed attacks) symmetrically.
